\subsection{CHEVALLEY ORDERS}

Let \((\mathfrak{g},\mathfrak{h})\) be a split reductive Lie algebra over \(\mathbf{Q}\) , \(\mathrm{R}\) its root system. Choose:

\begin{enumerate}
\def\labelenumi{\alph{enumi})}
\item
  a permissible lattice \(\mathcal{H}\) in \(\mathfrak{h}\) (no. 6, Def. 1);
\item
  for all \(\alpha \in \mathrm{R}\) , a lattice \(\mathcal{G}^{\alpha}\) in \(\mathfrak{g}^{\alpha}\) .
\end{enumerate}

Put \(\mathcal{G} = \mathcal{H} \oplus \sum_{\alpha \in \mathrm{R}} \mathcal{G}^{\alpha}\) . This is a lattice in \(\mathfrak{g}\) . Denote by \(\mathcal{U}\) the \(\mathbf{Z}\) -subalgebra of \(\mathrm{U}(\mathfrak{g})\) generated by the \(\binom{h}{n}\) ( \(h \in \mathcal{H}, n \in \mathbf{N}\) ) and the \(x^{(n)}\) ( \(x \in \mathcal{G}^{\alpha}, \alpha \in \mathrm{R}, n \in \mathbf{N}\) ). Finally, for \(\alpha \in \mathrm{R}\) and \(x \in \mathfrak{g}^{\alpha} - \{0\}\) , put

\[
w _ {\alpha} (x) = (\exp \operatorname{ad} x) (\exp \operatorname{ad} y) (\exp \operatorname{ad} x),
\]

where \(y\) is the unique element of \(\mathfrak{g}^{-\alpha}\) such that \([y,x] = H_{\alpha}\) . With these notations:

THEOREM 2. The following conditions are equivalent:

\begin{enumerate}
\def\labelenumi{(\roman{enumi})}
\item
  There exists a Chevalley system \((X_{\alpha})_{\alpha \in \mathrm{R}}\) of \((\mathfrak{g},\mathfrak{h})\) such that \(\mathcal{G}_{\alpha} = \mathbf{Z}X_{\alpha}\) for all \(\alpha \in \mathrm{R}\) .
\item
  \(\mathcal{U} \cap \mathfrak{g} = \mathcal{G}\) and \([\mathcal{G}^{\alpha}, \mathcal{G}^{-\alpha}] = \mathbf{Z}H_{\alpha}\) for all \(\alpha \in \mathrm{R}\) .
\item
  For all \(\alpha \in \mathrm{R}, x \in \mathcal{G}^{\alpha}, n \in \mathbf{N}\) , the endomorphism \((\operatorname{ad} x)^n / n!\) of \(\mathfrak{g}\) maps \(\mathcal{G}\) to \(\mathcal{G}\) , and \([\mathcal{G}^{\alpha}, \mathcal{G}^{-\alpha}] = \mathbf{Z}H_{\alpha}\) .
\item
  For all \(\alpha \in \mathrm{R}\) and every basis \(x\) of \(\mathcal{G}^{\alpha}\) , \(w_{\alpha}(x)\) maps \(\mathcal{G}\) to \(\mathcal{G}\) (that is, maps \(\mathcal{G}^{\beta}\) to \(\mathcal{G}^{s_{\alpha}(\beta)}\) for all \(\beta \in \mathrm{R}\) ).
\item
  \(\Longrightarrow\) (ii): let \((X_{\alpha})_{\alpha \in \mathrm{R}}\) be a Chevalley system in \((\mathfrak{g},\mathfrak{h})\) such that \(\mathcal{G}^{\alpha} = \mathbf{Z}X_{\alpha}\) for all \(\alpha \in \mathrm{R}\) , and let B be a basis of R. For \(\alpha \in \mathrm{B}\) , put \(x_{\alpha} = X_{\alpha},y_{\alpha} = X_{-\alpha}\) . Let \(\mathcal{U}'\) be the biorder associated by Prop. 3 of no. 6 to \(\mathcal{H}\) , the \(x_{\alpha}\) and the \(y_{\alpha}\) . It is clear that \(\mathcal{U}' \subset \mathcal{U}\) . By Cor. 3 and 4 of Prop. 3, \(x^{(n)} \in \mathcal{U}'\) for all \(\alpha \in \mathrm{R}\) , \(x \in \mathcal{G}^{\alpha}\) and \(n \in \mathbf{N}\) . Thus \(\mathcal{U} = \mathcal{U}'\) , which proves (ii).
\item
  \(\Longrightarrow\) (iii): this is clear by Lemma 2 of no. 5.
\item
  \(\Longrightarrow\) (iv): let \(\alpha \in \mathrm{R}\) and let \(x\) be a basis of \(\mathcal{G}^{\alpha}\) . Since \([\mathcal{G}^{\alpha},\mathcal{G}^{-\alpha}] = \mathbf{Z}H_{\alpha}\) , the unique \(y \in \mathfrak{g}^{-\alpha}\) such that \([y,x] = H_{\alpha}\) belongs to \(\mathcal{G}^{-\alpha}\) . Since \(\exp \mathrm{ad} x\) and \(\exp \mathrm{ad} y\) leave \(\mathcal{G}\) stable by (iii), so does \(w_{\alpha}(x)\) .
\item
  \(\Longrightarrow\) (i): let B be a basis of R. Choose a basis \(x_{\alpha}\) of \(\mathcal{G}^{\alpha}\) for all \(\alpha \in \mathrm{B}\) . Let \(y_{\alpha} \in \mathcal{G}^{-\alpha}\) be such that \([y_{\alpha}, x_{\alpha}] = H_{\alpha}\) . By §1, no. 5, formulas (5), we have \(y_{\alpha} = w_{\alpha}(x_{\alpha}). x_{\alpha}\) so \(y_{\alpha}\) is a basis of \(\mathcal{G}^{-\alpha}\) by (iv). Let \(\mathcal{G}\) be the order in \(\mathfrak{g}\) defined by \(\mathcal{H}\) , the \(x_{\alpha}\) and the \(y_{\alpha}\) (no. 6, Cor. 1 of Prop. 3). Then \(\mathcal{G}\) is stable under the \((\operatorname{ad} x_{\alpha})^{n} / n!\) , \((\operatorname{ad} y_{\alpha})^{n} / n!\) (loc. cit.), and hence under the \(w_{\alpha}(x_{\alpha})\) .
\end{enumerate}

Now let \(\beta \in \mathrm{R}\) . There exist \(\alpha_0, \alpha_1, \ldots, \alpha_r \in \mathrm{B}\) such that

\[
\beta = s _ {\alpha_ {r}} s _ {\alpha_ {r - 1}} \dots s _ {\alpha_ {1}} (\alpha_ {0})
\]

(Chap. VI, §1, no. 5, Prop. 15). Then \(w_{\alpha_r}(x_{\alpha_r}).w_{\alpha_{r-1}}(x_{\alpha_{r-1}}) \ldots w_{\alpha_1}(x_{\alpha_1})\) maps \(\mathcal{G}^{\alpha_0}\) to \(\mathcal{G}^\beta\) by (iv), and maps \(\mathcal{G} \cap \mathfrak{g}^{\alpha_0}\) to \(\mathcal{G} \cap \mathfrak{g}^\beta\) by the preceding. Since \(\mathcal{G} \cap \mathfrak{g}^{\alpha_0} = \mathcal{G}^{\alpha_0}\) (Prop. 3 (ii)), we have \(\mathcal{G} \cap \mathfrak{g}^\beta = \mathcal{G}^\beta\) . Thus

\[
\mathcal {G} = \mathcal {H} \oplus \sum_ {\beta \in \mathrm{R}} (\mathcal {G} \cap \mathfrak {g} ^ {\beta}) = \mathcal {H} \oplus \sum_ {\beta \in \mathrm{R}} \mathcal {G} ^ {\beta} = \mathcal {G}
\]

and Cor. 4 of Prop. 3 concludes the proof.

DEFINITION 2. When conditions (i) to (iv) of Th. 2 are satisfied, \(\mathcal{G}\) is said to be a Chevalley order in \((\mathfrak{g}, \mathfrak{h})\) .

Remark. Chevalley orders in \((\mathfrak{g},\mathfrak{h})\) always exist. Indeed, the Chevalley orders are the sets of the form \(\mathcal{H}\oplus\sum_{\alpha\in \mathrm{R}}\mathbf{Z}X_{\alpha}\) , where \((X_{\alpha})_{\alpha\in \mathrm{R}}\) is a Chevalley system in \((\mathfrak{g},\mathfrak{h})\) and \(\mathcal{H}\) is a lattice in \(\mathfrak{h}\) such that

\[
\mathrm{Q} (\mathrm{R} ^ {\vee}) \subset \mathcal {H} \subset \mathrm{P} (\mathrm{R} ^ {\vee}) \oplus \mathfrak {c}
\]

(c being the centre of g).

THEOREM 3. We retain the notations at the beginning of no. 7, and assume that \(\mathcal{G}\) is a Chevalley order in \((\mathfrak{g},\mathfrak{h})\) .

\begin{enumerate}
\def\labelenumi{(\roman{enumi})}
\item
  \(\mathcal{U}\) is a b-order in \(\mathrm{U}(\mathfrak{g})\) .
\item
  Let B be a basis of \(\mathrm{R}\), and \((X_{\alpha})_{\alpha\in\mathrm{B}\cup(-\mathrm{B})}\) a family of elements of \(\mathfrak{g}\) such that \(\mathcal{G}^{\alpha}=\mathbf{Z}X_{\alpha}\) for \(\alpha\in\mathrm{B}\cup(-\mathrm{B})\) . The \(\mathbf{Z}\)-algebra \(\mathcal{U}\) is generated by the \(\binom{h}{n}\) and the \(X_{\alpha}^{(n)}\) ( \(h\in\mathcal{H},\alpha\in\mathrm{B}\cup(-\mathrm{B}),n\in\mathbf{N}\) ). If \(\mathfrak{g}\) is semi-simple and \(\mathcal{H}=\mathrm{Q}(\mathrm{R}^{\vee})\) , the \(\mathbf{Z}\)-algebra \(\mathcal{U}\) is generated by the \(X_{\alpha}^{(n)}\) ( \(\alpha\in\mathrm{B}\cup(-\mathrm{B}),n\in\mathbf{N}\) ).
\item
  Let B be a basis of R, \(R_{+}\) the corresponding set of positive roots, \(R_{-} = -R_{+}\) , \(n_{+} = \sum_{\alpha \in R_{+}} g^{\alpha}\) , \(n_{-} = \sum_{\alpha \in R_{-}} g^{\alpha}\) . Then,
\end{enumerate}

\[
\mathcal {U} = (\mathcal {U} \cap \mathrm{U} (\mathfrak {n} _ {-})). (\mathcal {U} \cap \mathrm{U} (\mathfrak {h})). (\mathcal {U} \cap \mathrm{U} (\mathfrak {n} _ {+})).
\]

Let \((h_i)_{i\in \mathrm{I}}\) be a basis of \(\mathcal{H}\) . For all \(\alpha \in \mathrm{R}\) , let \(X_{\alpha}\) be a basis of \(\mathcal{G}^{\alpha}\) . Give the set \(\mathrm{I}\cup \mathrm{R}\) a total order (we assume that \(\mathrm{I}\cap \mathrm{R} = \emptyset\) ). For \(\lambda \in \mathrm{I}\cup \mathrm{R}\) and \(n\in \mathbf{N}\) , put \(e_{\lambda}^{\langle n\rangle} = \binom{h_{\lambda}}{n}\) if \(\lambda \in \mathrm{I}\) , \(e_{\lambda}^{\langle n\rangle} = X_{\lambda}^{(n)}\) if \(\lambda \in \mathrm{R}\) . Then the products \(\prod_{\lambda \in \mathrm{I}\cup \mathrm{R}}e_{\lambda}^{\langle n_{\lambda}\rangle}\) , where \((n_{\lambda})\) belongs to \(\mathbf{N}^{\mathrm{I}\cup \mathrm{R}}\) , form a basis of the \(\mathbf{Z}\) -module \(\mathcal{U}\) . The products \(\prod_{\lambda \in \mathrm{I}}\binom{h_{\lambda}}{n_{\lambda}}\) , where \((n_{\lambda})\) belongs to \(\mathbf{N}^{\mathrm{I}}\) , form a basis of the \(\mathbf{Z}\) -module \(\mathcal{U}\cap \mathrm{U}(\mathfrak{h})\) . The products \(\prod_{\lambda \in \mathrm{R}_{+}}X_{\lambda}^{(n_{\lambda})}\) , where \((n_{\lambda})\) belongs to \(\mathbf{N}^{\mathrm{R}_{+}}\) , form a basis of the \(\mathbf{Z}\) -module \(\mathcal{U}\cap \mathrm{U}(\mathfrak{n}_{+})\) .

Let B and \((X_{\alpha})_{\alpha\in\mathrm{B}\cup(-\mathrm{B})}\) be as in (ii), and such that \([X_{-\alpha},X_{\alpha}]=H_{\alpha}\) . Let \(\mathcal{U}'\) be the \(\mathbf{Z}\)-subalgebra of \(\mathrm{U}(\mathfrak{g})\) generated by the \(\binom{h}{n}\) and the \(X_{\alpha}^{(n)}\) \((h\in\mathcal{H},\alpha\in\mathrm{B}\cup(-\mathrm{B}),n\in\mathbf{N})\) . We have seen in the proof of Th. 2, (i) \(\Longrightarrow\) (ii), that \(\mathcal{U}'\) is equal to \(\mathcal{U}\) and is a biorder in \(\mathrm{U}(\mathfrak{g})\) . This proves (i) and the first assertion of (ii); the second follows from Lemma 4 (ii). Assertion (iii) follows from Th. 1 (no. 3) and Prop. 3 (no. 6).

